\begin{textsample}{POS Dim 3 – es – Score 16.00 – es/es\_inf\_5\_acolhimento.txt}  \label{ex:f3_pos_065}
ACOLHIMENTO Acolhimento na Educação \textbf{Infantil} significa mais que cumprimentar e receber no portão ou nos espaços da \textbf{creche}, o \textbf{bebê}, a \textbf{criança} e sua família ou despedir -se deles à saída.
Fazem parte do acolhimento a sensibilidade, a afetividade, o \textbf{cuidado}, a delicadeza no trato, a cordialidade, a \textbf{escuta} qualificada às \textbf{crianças} e aos seus pais.
Significa estabelecer uma interlocução respeitosa, humanizada, identificando e chamando as \textbf{crianças} e seus familiares pelos seus nomes, com uma atitude de quem se importa, de quem respeita sua cultura, valoriza suas histórias de vida, experiências e vivências.
O acolhimento faz gerar uma relação de pertencimento dos \textbf{bebês} e das \textbf{crianças} com o espaço escolar, criando um vínculo afetivo e de segurança na relação com os profissionais da escola.
Viabiliza canais de diálogo, estabelecendo uma relação de empatia, de confiança e de respeito com a \textbf{criança} e com seus familiares.
Acolhida pode traduzir -se numa oportunidade de a \textbf{criança} receber dos educadores um tempo de atenção, de carinho, de \textbf{brincadeiras} e de \textbf{afeto} que favorecem o seu desenvolvimento intra e interpessoal.
A postura acolhedora deve caracterizar uma prática constante dos profissionais da educação, podendo acontece por meio das \textbf{pequenas} e corriqueiras ações do dia a dia, que pode ser um \textbf{colo} ou uma canção.
Os/as professores/as e demais profissionais que interagem com as \textbf{crianças} nos momentos de alimentação, de trocas de \textbf{fraldas}, de banho, de \textbf{brincadeiras}, de vivências individuais e nos grupos devem fazer dessas ações momentos para \textbf{acolher} e valorizar os sentimentos da \textbf{criança}, dando -lhe oportunidade para manifestar suas emoções, expressões, linguagens, comunicando -lhes \textbf{afeto}, estimulando suas interações e aprendizagem, visto que o educar e o cuidar são ações indissociáveis na educação \textbf{infantil}.
O acolhimento aos familiares das \textbf{crianças} define o tom da relação que estes terão com a instituição.
O documento Práticas Cotidianas na Educação \textbf{Infantil} ( Brasil, 2009, p. 33 ) esclarece que : > acontecem na instituição educativa, afinal as \textbf{crianças} são \textbf{pequenas} e para se sentirem \textbf{acolhidas} na \textbf{creche}, dependem da sintonia entre a família e os profissionais da escola.
Essa é uma das especificidades dos estabelecimentos de Educação \textbf{Infantil}.
Nesse sentido, complementaridade e partilha são palavras decisivas na relação escola, \textbf{criança} e família.¹ RELAÇÃO FAMÍLIA-ESCOLA Consideramos o conceito família no seu sentido amplo e inclusivo : um grupo de duas ou mais pessoas, não necessariamente ligadas por laços de consanguinidade, mas unidas por laços afetivos e por compromissos recíprocos.
As instituições de Educação \textbf{Infantil} devem estar preparadas para lidar com as diferentes estruturas familiares, considerando legítima a participação, não apenas da família natural, mas da substituta, da de guarda e tutela, de todas as que exercem funções insubstituíveis de proteção, de assistência e \textbf{cuidados}, de educação e promoção de valores.
Todas devem ter garantidos e respeitados seus direitos de participação nos processos de educação e \textbf{cuidados} das \textbf{crianças}.
A Convenção sobre os Direitos da Criança, Organização das Nações Unidas, postula que : > Para um desenvolvimento completo e harmonioso de sua personalidade, a \textbf{criança} deve crescer num ambiente familiar, numa atmosfera de felicidade, amor e compreensão.
Portanto, educar as \textbf{crianças} é uma tarefa que exige tempo, amor, \textbf{cuidados} e conhecimentos tanto por parte da família quanto das instituições que as recebem.
Quando as famílias e as instituições educativas têm identidade de propósitos e objetivos comuns, com relação aos \textbf{cuidados} e à educação das \textbf{crianças}, as condições de desenvolvimento pessoal e social destas são favorecidas pelo diálogo e se ampliam.
Na tarefa de educar, seja no lar ou na escola, o \textbf{afeto} tem papel preponderante, pois deixa marcas no desenvolvimento psíquico da \textbf{criança}.
Portanto, todo o aprendizado dos \textbf{bebês} e das \textbf{crianças}, seja na família ou nas instituições, deve ser conduzido e mediado pelo \textbf{afeto}.
Ao estabelecer relações com as famílias a instituição de educação deve levar em conta que estas têm histórias, culturas próprias, que trazem as marcas das relações e experiências dos seus antepassados, o que as tornam diversas e singulares.
O \textbf{bebê} nasce nesse ambiente familiar, no qual as primeiras interações sociais acontecem, constituindo -se no espaço para o desenvolvimento da sua consciência de ser humano.
Ainda \textbf{pequeno}, ele é capaz de reagir ao seu entorno e produzir reações nos que o cercam e, como ser em crescimento, ele necessita de estímulos múltiplos, que potencializem seu desenvolvimento, tanto no plano físico quanto no psicológico. ¹ Convenção sobre os Direitos da Criança, Organização das Nações Unidas, 1989 A formação plena da \textbf{criança}, de acordo com os princípios legais, requer esforços integrados, compromissos e compartilhamento de responsabilidade entre famílias, instituições de educação e a sociedade, a fim de assegurar que seus direitos sejam respeitados.
Ao ingressarem em contextos coletivos da Educação \textbf{Infantil}, os \textbf{bebês} e as \textbf{crianças} têm o seu universo pessoal de significados ampliado e nas interações sociais com \textbf{adultos} e \textbf{crianças} maiores vão, \textbf{progressivamente}, construindo características identitárias, qualidades de caráter e atitudes, que se revelam cotidianamente em suas ações e relações.
Contudo, a convivência das \textbf{crianças} com toda essa diversidade, nos espaços da Educação \textbf{Infantil}, torna -se mais rica quando a família participa.
Assim, família e escola devem comungar dos mesmos objetivos e propostas de formação integral da \textbf{criança}, que consistem do desenvolvimento cognitivo, da imaginação e das competências socioemocionais, que determinam a constituição da sua identidade e autonomia.
Para participar a família precisa conhecer a proposta de educação da escola para suas \textbf{crianças}.
Às vezes é necessário que a instituição mostre aos pais os caminhos para a sua participação, pois embora tenham interesse, não sabem como, sendo necessário que a instituição lhes esclareça sobre seus direitos e deveres de participação.
Assim, de acordo com o estabelece a Lei 9394/96 de Diretrizes e Bases da Educação Nacional – LDB : > Art. 1º A educação abrange os processos formativos que se desenvolvem na vida familiar, na convivência humana, no trabalho, nas instituições de ensino e pesquisa, nos movimentos sociais e organizações da sociedade civil e nas manifestações culturais.
A troca de conhecimentos entre as famílias e os profissionais da escola, sobre os processos de educação, valores e expectativas, o \textbf{acompanhamento} das vivências cotidianas das \textbf{crianças} pelos pais ou responsáveis auxiliam no desenvolvimento, na inserção e adaptação destas aos ambientes da \textbf{creche} e da \textbf{pré-escola}, e influenciam na constituição da sua autoestima e no seu desenvolvimento.
Portanto, família e escola devem estar juntas nessa grande empreitada de \textbf{apoiar} e estimular as \textbf{crianças} nas suas vivências, na descoberta de suas potencialidades, dos seus gostos, das suas dificuldades, como parceiras nos processos de cuidar e educar.

% matched lemmas: acolher, acompanhamento, adulto, afeto, apoiar, bebê, brincadeira, colo, creche, criança, cuidado, escuta, fralda, infantil, pequeno, progressivamente, pré-escola
\end{textsample}
